\scanpage{67}
\mspage{33}
``forage de trous'' doit nous fournir un groupe extérieur $\Pi_{0,3}(I)$, \inserted{de type $(0,3+\operatorname{card} I)$,} sur lequel $\Gamma$ opère (il vaut mieux peut-être utiliser le yoga introduit par ailleurs des groupoïdes rigides — donc on \unsure{peut|pourra} \unsure{ou devra} aussi boucher certains des trous $0,1,\infty$ — on trouve donc l'équivalent groupoïdal de la \deleted{``sphère \unsure{ar.}} droite projective $\mathbb P^1_Q$, on l'appelle la sphère arithmétique — qui correspond à un groupe \inserted{ext.} à lacets infini sur lequel $\Gamma$ opère — en fait, ce n'est autre que $E_{K_1}$, où $\unsure{\vert}$
\[
\tag{6}K_1=Q(\deleted{\unsure{X_1|x_1}}\,T_1)
\]
est l'extension transcendante pure \inserted{typ.} de degré $1$ de $Q$.

Partant de (6), on construit de $\unsure{\overline m}$ l'équivalent groupoïdal de $\deleted{\unsure{X|P}}\,U_{0,3}$ $\unsure{\longleftarrow}$ et on remonte comme
\margin{66}
